\section*{Notes and Reading}
\addcontentsline{toc}{section}{Notes and Reading}

The references below make the empirical claims in this chapter auditable and distinguish published peer-reviewed findings from interpretive strategic models. Historical U.S. survival and earnings baselines should not be treated as deterministic forecasts for an individual venture or an Indian market context. All numerical models presented are explicitly illustrative.

\begin{enumerate}[label={[\arabic*]},leftmargin=14pt,itemsep=3pt,topsep=2pt,parsep=0pt]
  \item Paul Graham, ``How to Make Wealth'' (May 2004). Distinguishes money from wealth and investigates how small teams create and capture economic surplus. Available at: \href{https://www.paulgraham.com/wealth.html}{paulgraham.com/wealth.html}.
  \item Naval Ravikant, ``How to Get Rich (without getting lucky)'' (2019--2020). Examines specific knowledge, accountability, and permissioned versus permissionless leverage. Available at: \href{https://nav.al/rich}{nav.al/rich}.
  \item Paul Graham, ``Do Things that Don't Scale'' (July 2013). Focuses on manual customer acquisition and intensive qualitative feedback during beachhead formation. Available at: \href{https://paulgraham.com/ds.html}{paulgraham.com/ds.html}.
  \item Paul Graham, ``Startup = Growth'' (September 2012). Defines the essential structure of a startup as an enterprise designed for rapid scale across a large market. Available at: \href{https://www.paulgraham.com/growth.html}{paulgraham.com/growth.html}.
  \item Erik Hurst and Benjamin Wild Pugsley, ``What Do Small Businesses Do?'' \textit{NBER Working Paper} No. 17041 (2011). Empirical analysis of nonfinancial motivations, growth aspirations, and lifestyle choices among small-business owners. Available at: \href{https://www.nber.org/papers/w17041}{nber.org/papers/w17041}.
  \item Warren Buffett, Berkshire Hathaway Shareholder Letter (1986). Seminal treatment of economic goodwill, capital allocation, and owner earnings versus reported accounting profit. Available at: \href{https://www.berkshirehathaway.com/letters/1986.html}{berkshirehathaway.com/letters/1986.html}.
  \item Barton H. Hamilton, ``Does Entrepreneurship Pay? An Empirical Analysis of the Returns to Self-Employment,'' \textit{Journal of Political Economy} 108(3), 604--631 (2000). Documents a substantial median earnings disadvantage in self-employment relative to paid employment. Available at: \href{https://doi.org/10.1086/262131}{doi.org/10.1086/262131}.
  \item Tobias J. Moskowitz and Annette Vissing-J\o{}rgensen, ``The Returns to Entrepreneurial Investment: A Private Equity Premium Puzzle?'' \textit{American Economic Review} 92(4), 745--778 (2002). Analyzes the lack of an aggregate risk premium for holding concentrated private equity. Available at: \href{https://www.nber.org/papers/w8876}{nber.org/papers/w8876}.
  \item Robert E. Hall and Susan E. Woodward, ``The Burden of the Nondiversifiable Risk of Entrepreneurship,'' \textit{American Economic Review} 100(3), 1163--1194 (2010). Quantifies the distribution of net cash outcomes for venture-backed founders under exit uncertainty. Available at: \href{https://www.nber.org/papers/w14219}{nber.org/papers/w14219}.
  \item U.S. Bureau of Labor Statistics, ``Business Employment Dynamics: Twentieth Anniversary!'' (2024). Multi-cohort establishment survival data tracking firm longevity across longitudinal census cycles. Available at: \href{https://www.bls.gov/spotlight/2024/business-employment-dynamics-twentieth-anniversary/}{bls.gov/spotlight/2024/BED-report}.
  \item Pierre Azoulay, Benjamin Jones, J. Daniel Kim, and Javier Miranda, ``Age and High-Growth Entrepreneurship,'' \textit{American Economic Review: Insights} 2(1), 65--82 (2020). Demonstrates that the mean age of founders of the fastest-growing 0.1\% of new ventures is 45, driven by domain expertise and managerial capital. Available at: \href{https://www.nber.org/papers/w24489}{nber.org/papers/w24489}.
\end{enumerate}
